\documentclass[11pt]{article}
\usepackage[letterpaper,margin=0.72in]{geometry}
\usepackage[T1]{fontenc}
\usepackage[utf8]{inputenc}
\usepackage{lmodern}
\usepackage{microtype}
\usepackage{xcolor}
\usepackage{graphicx}
\usepackage{booktabs}
\usepackage{longtable}
\usepackage{tabularx}
\usepackage{array}
\usepackage{enumitem}
\usepackage[most]{tcolorbox}
\usepackage{fancyhdr}
\usepackage{titlesec}
\usepackage{hyperref}
\usepackage{xurl}
\usepackage{lastpage}
\usepackage{ragged2e}
\usepackage{multicol}
\usepackage{pdflscape}
\usepackage{amssymb}
\definecolor{BAOrange}{HTML}{E87521}
\definecolor{BANavy}{HTML}{17324D}
\definecolor{BASlate}{HTML}{5B6573}
\definecolor{BALight}{HTML}{F4F6F8}
\definecolor{BAGold}{HTML}{D7A62A}
\hypersetup{colorlinks=true,linkcolor=BANavy,urlcolor=BAOrange,citecolor=BANavy,pdftitle={Engineering and Architecture Crosswalk}}
\setlength{\parindent}{0pt}
\setlength{\parskip}{4pt}
\setlist{nosep,leftmargin=*}
\renewcommand{\arraystretch}{1.18}
\newcolumntype{Y}{>{\RaggedRight\arraybackslash}X}
\newcolumntype{L}[1]{>{\RaggedRight\arraybackslash}p{#1}}
\pagestyle{fancy}
\setlength{\headheight}{14pt}
\fancyhf{}
\lhead{\textcolor{BANavy}{\small BABOK v3 Study Series}}
\rhead{\textcolor{BASlate}{\small Engineering and Architecture Crosswalk}}
\cfoot{\small \thepage\ / \pageref{LastPage}}
\renewcommand{\headrulewidth}{0.4pt}
\titleformat{\section}{\Large\bfseries\color{BANavy}}{\thesection}{0.7em}{}
\titleformat{\subsection}{\large\bfseries\color{BAOrange}}{\thesubsection}{0.6em}{}
\titleformat{\subsubsection}{\normalsize\bfseries\color{BANavy}}{\thesubsubsection}{0.5em}{}
\newtcolorbox{keybox}[1][]{enhanced,breakable,colback=BALight,colframe=BANavy,boxrule=0.7pt,arc=1.5mm,left=2mm,right=2mm,top=1.5mm,bottom=1.5mm,#1}
\newtcolorbox{practicebox}[1][]{enhanced,breakable,colback=white,colframe=BAOrange,boxrule=0.8pt,arc=1.5mm,left=2mm,right=2mm,top=1.5mm,bottom=1.5mm,#1}
\newtcolorbox{warningbox}[1][]{enhanced,breakable,colback=white,colframe=BAGold,boxrule=0.8pt,arc=1.5mm,left=2mm,right=2mm,top=1.5mm,bottom=1.5mm,#1}
\begin{document}
\begin{titlepage}
\centering
\vspace*{1.15in}
{\Huge\bfseries\color{BANavy} Engineering and Architecture Crosswalk\par}
\vspace{0.35in}
{\Large\color{BAOrange} BABOK practices mapped to software engineering, AppSec, Agile, UML/BPMN, and enterprise architecture\par}
\vspace{0.55in}
\rule{0.78\textwidth}{1.2pt}\par
\vspace{0.35in}
{\large Technical reference and study artifact based on the uploaded \emph{A Guide to the Business Analysis Body of Knowledge (BABOK Guide), Version 3}.\par}
\vfill
\begin{keybox}
\textbf{Source discipline.} BABOK terminology and organization are preserved where practical. Explanations are paraphrased for study and engineering use. Any engineering, security, Agile, UML/BPMN, or enterprise-architecture crosswalks are practitioner mappings, not normative BABOK statements.
\end{keybox}
\vspace{0.35in}
{\small Prepared October 2026\par}
\end{titlepage}
\tableofcontents
\newpage
\section{Interpretation Boundary}
The BABOK source explicitly contains Agile, Information Technology, Business Architecture, and Business Process Management perspectives, and it includes modelling techniques such as process modelling, sequence diagrams, state modelling, use cases/scenarios, data modelling, and roles/permissions matrices. The detailed crosswalk below extends those practices into modern software engineering and application security. Those extensions are practitioner mappings, not claims that BABOK prescribes a particular SDLC, security framework, UML diagram set, or enterprise-architecture method.
\section{Domain Lens Summary}
\begin{longtable}{L{0.18\textwidth} L{0.70\textwidth}}
\toprule
Software Engineering & Treat requirements as part of a traceable engineering system spanning architecture, code, tests, interfaces, releases, operations, and feedback.\\
Application Security & Add assets, threats, abuse cases, security quality requirements, control design, verification evidence, risk acceptance, and remediation outcomes to the same value/traceability model.\\
Agile & Preserve outcomes and traceability while reducing speculative detail. Analyze continuously, refine just in time, prioritize by value/risk, and use increments to learn.\\
UML / BPMN & Select notation based on the question: actors/goals, interactions, states, structure, or business workflow. Models are evidence and communication artifacts, not substitutes for decisions or traceability.\\
Enterprise Architecture & Connect needs and requirements to capabilities, operating model, target state, architecture requirements, transition states, roadmaps, standards, and governance.\\
\bottomrule
\end{longtable}
\section{Thirty-Task Crosswalk}
\begin{landscape}
\footnotesize
\setlength{\tabcolsep}{3pt}
\begin{longtable}{L{0.11\linewidth} L{0.16\linewidth} L{0.16\linewidth} L{0.15\linewidth} L{0.16\linewidth} L{0.16\linewidth}}
\toprule
\textbf{BABOK task} & \textbf{Software engineering} & \textbf{AppSec} & \textbf{Agile} & \textbf{UML / BPMN} & \textbf{Enterprise architecture}\\
\midrule
\endfirsthead
\toprule
\textbf{BABOK task} & \textbf{Software engineering} & \textbf{AppSec} & \textbf{Agile} & \textbf{UML / BPMN} & \textbf{Enterprise architecture}\\
\midrule
\endhead
3.1 Plan BA Approach & SDLC/analysis plan; define artifacts and review cadence & Include security analysis activities, evidence, and gates in the plan & Tailor analysis to product cadence; just-in-time detail & Select required model types and modelling conventions & Align analysis method with architecture governance and lifecycle\\
3.2 Plan Stakeholder Engagement & Identify product, engineering, operations, QA, compliance, and business participants & Identify security owners, security champions, risk owners, privacy, IAM, SOC/IR as relevant & Clarify Product Owner, team, users, SMEs, and decision makers & Identify actors, lanes, participants, model reviewers & Capture stakeholder concerns that architecture views must address\\
3.3 Plan BA Governance & Define decision rights, reviews, baselines, and change control & Define security risk acceptance, exception/waiver, severity and sign-off paths & Embed policy in Definition of Ready/Done and backlog governance & Use model reviews as decision evidence; preserve approved baselines & Connect requirements/design decisions to architecture review and ADR governance\\
3.4 Plan BA Information Management & Define IDs, repositories, versioning, metadata, and links & Store security requirements, threats, findings, exceptions, evidence, and ownership & Link epics/stories/acceptance criteria to supporting artifacts & Define model repository, naming, versioning, and trace links & Integrate with architecture repository, catalog, capability/process models\\
3.5 Identify BA Performance Improvements & Measure requirement quality, rework, review cycle time, defect leakage & Measure security requirement coverage, false-positive burden, remediation latency, escape rate & Use retrospectives and flow metrics to improve analysis & Track model usefulness, defects, stale artifacts & Feed lessons into architecture governance and standards\\
4.1 Prepare for Elicitation & Plan discovery interviews, workshops, observation, data review & Prepare threat/abuse-case elicitation, security questionnaires, regulatory inputs & Prepare refinement/discovery sessions near delivery need & Select use case, process, state, sequence, context modelling sessions & Prepare capability, value-stream, stakeholder, and operating-model discovery\\
4.2 Conduct Elicitation & Discover business behavior, constraints, interfaces, data, quality attributes & Elicit assets, trust boundaries, abuse cases, security objectives, compliance needs & Use collaborative refinement and continuous discovery & Build use cases, BPMN flows, context/sequence/state models with SMEs & Elicit strategy, capability gaps, principles, dependencies, transition constraints\\
4.3 Confirm Elicitation Results & Reconcile notes, models, acceptance criteria, assumptions & Confirm security assumptions, attack surfaces, data classifications, control intent & Confirm stories/criteria with team and stakeholders before commitment & Walk through models against scenarios and edge cases & Validate enterprise facts and concerns with domain owners\\
4.4 Communicate BA Information & Publish audience-specific requirements, diagrams, decisions & Translate technical risk and control needs for engineers and leaders & Use concise backlog items, examples, demos, lightweight models & Use the notation best suited to the question; avoid model overload & Present architecture concerns through stakeholder-relevant views\\
4.5 Manage Stakeholder Collaboration & Resolve conflicts and sustain participation across delivery roles & Coordinate engineering, security, risk, compliance, and operations decisions & Facilitate team collaboration and rapid feedback & Use shared modelling to expose disagreements explicitly & Maintain alignment across business, solution, data, security, and technology architecture\\
5.1 Trace Requirements & Link objectives -> requirements -> design -> code/change -> test -> release & Link security objectives/threats -> controls/requirements -> implementation -> tests/scans -> findings/acceptance & Link epics/features/stories/criteria to tests and releases without excessive documentation & Use diagrams as trace evidence; keep authoritative links in ALM/requirements tooling & Trace business outcomes and capabilities to architecture building blocks and initiatives\\
5.2 Maintain Requirements & Keep baselines current as code, interfaces, and behavior evolve & Maintain security requirements as threats, dependencies, and controls change & Continuously refine backlog; retire obsolete items & Version models with the behavior they describe & Maintain architecture requirements, principles, standards, and roadmaps\\
5.3 Prioritize Requirements & Order work using value, risk, dependencies, cost, and urgency & Include exploitability, exposure, control gaps, regulatory deadlines, and risk reduction & Backlog ordering and iteration/release planning & Use model dependencies to reveal sequencing constraints & Prioritize capability gaps and architecture work by strategic value and dependency\\
5.4 Assess Requirements Changes & Perform impact analysis across design, code, tests, interfaces, schedule & Reassess threat model, security controls, attack surface, risk, and required evidence & Evaluate change against committed iteration and future backlog & Trace changed flows/states/interfaces through affected models & Assess effects on target architecture, standards, transition states, and roadmaps\\
5.5 Approve Requirements & Record agreement needed for build/test/release & Obtain required security/risk acceptance without treating approval as proof of security & Use team/product decision policies appropriate to the increment & Record model/specification reviews where governance requires & Obtain architecture/governance decisions for material requirements or exceptions\\
6.1 Analyze Current State & Baseline systems, processes, interfaces, pain points, constraints, and quality & Assess current attack surface, vulnerabilities, controls, data flows, security debt & Understand current product behavior and operational feedback & Model as-is process/system interactions and key states & Baseline capabilities, value streams, organization, applications, data, technology\\
6.2 Define Future State & Define target outcomes, quality attributes, capabilities, and operating conditions & Define target security posture, trust model, security outcomes, resilience objectives & Express product goals and measurable outcomes before decomposing backlog & Model to-be flows, interactions, states, and boundaries & Define target capabilities and target architecture outcomes\\
6.3 Assess Risks & Evaluate delivery, operational, technical, dependency, and adoption risk & Perform threat/risk analysis; include likelihood, impact, exposure, compensating controls & Prioritize uncertainty and spikes; make risk visible in backlog & Use misuse/abuse scenarios and model boundary analysis to expose risk & Assess architecture, transition, vendor, data, technology, and organizational risk\\
6.4 Define Change Strategy & Choose migration/release approach and transition states & Plan control rollout, remediation strategy, compensating controls, and security verification & Use incremental releases and feedback to reduce uncertainty & Model transition processes and interim interfaces/states & Create transition architectures, work packages, dependencies, and roadmap\\
7.1 Specify and Model Requirements & Define functional, interface, data, and quality requirements in testable form & Define security requirements, abuse cases, security NFRs, logging/audit, IAM, data protection & Express backlog items with examples and acceptance criteria at needed detail & Use UML use case/sequence/state/class where useful; BPMN for business workflows & Define architecture requirements and viewpoints needed to constrain solution design\\
7.2 Verify Requirements & Check requirement/model quality, consistency, ambiguity, completeness, and testability & Check security requirements are specific, measurable/testable, scoped, and internally consistent & Use peer review/refinement and executable examples where possible & Validate notation syntax/semantics and consistency across diagrams & Review architecture requirements for clarity, consistency, traceability, and feasibility\\
7.3 Validate Requirements & Confirm requirements support the actual business need and desired outcomes & Confirm controls reduce relevant risk and do not merely satisfy a checklist & Validate backlog value with users/product outcomes & Walk scenarios through models to confirm they solve the intended problem & Confirm architecture requirements support stakeholder concerns and target outcomes\\
7.4 Define Requirements Architecture & Organize requirements by layers, domains, features, components, dependencies & Organize security requirements by asset, trust zone, control family, threat, component & Structure epics/features/stories while preserving outcome links & Use model packages/views and explicit relationships; do not rely on diagrams alone for traceability & Structure requirements around capabilities, domains, architecture layers, and viewpoints\\
7.5 Define Design Options & Compare feasible solution approaches and allocate requirements & Compare security architectures and control patterns; include residual risk and operational burden & Use spikes/prototypes to compare options before commitment & Represent alternatives with sequence, state, component/context, and process models & Compare target-state alternatives against principles, standards, capabilities, and roadmap constraints\\
7.6 Analyze Potential Value and Recommend Solution & Compare options by value, cost, risk, feasibility, maintainability, and fit & Include risk reduction, attack surface, assurance cost, security operability, and residual risk & Choose the smallest valuable solution that supports learning and delivery goals & Use models to explain trade-offs rather than as the decision by themselves & Recommend architecture option based on strategic fit, value, risk, and transition feasibility\\
8.1 Measure Solution Performance & Define product/system measures tied to expected outcomes & Measure vulnerability trends, control effectiveness, exploitability, remediation, coverage, security events & Use product and flow metrics plus operational telemetry & Instrument measures represented by important states/process outcomes & Measure capability and architecture outcomes, not only project completion\\
8.2 Analyze Performance Measures & Explain whether observed results indicate value and where gaps exist & Interpret security telemetry and findings in context; separate volume from risk & Review outcomes in demos/reviews/retrospectives and adjust backlog & Use process/state models to localize performance bottlenecks or failure paths & Assess whether architecture changes improved strategic outcomes\\
8.3 Assess Solution Limitations & Find design, code, configuration, interface, usability, or operational weaknesses & Analyze vulnerabilities, insecure design, missing controls, misconfiguration, weak observability & Turn evidence into improvement stories/defects & Use interaction/state/process models to isolate failure mechanisms & Identify solution architecture constraints and technical debt limiting value\\
8.4 Assess Enterprise Limitations & Find organizational, process, policy, skill, data, funding, or governance barriers & Identify security process, ownership, tooling, skills, incentives, policy, and governance barriers & Surface systemic impediments beyond a single team & Use BPMN/organization models to expose handoff and accountability problems & Identify capability gaps, operating-model issues, or governance constraints\\
8.5 Recommend Actions to Increase Solution Value & Recommend optimization, enhancement, workaround, replacement, or retirement & Recommend prioritized remediation, compensating controls, architecture changes, or retirement & Feed validated improvements into backlog and roadmap & Use before/after models to communicate the change & Update roadmaps, transition architectures, standards, and capability investments\\
\bottomrule
\end{longtable}
\normalsize
\end{landscape}
\section{Model Selection: UML and BPMN}
\begin{longtable}{L{0.20\textwidth} L{0.32\textwidth} L{0.36\textwidth}}
\toprule
\textbf{Question} & \textbf{Useful representation} & \textbf{BABOK connection}\\
\midrule
Who interacts with the solution to achieve a goal? & UML use-case diagram plus textual scenarios & Use Cases and Scenarios; Stakeholder List/Map/Personas; Specify and Model Requirements.\\
How do participants exchange messages over time? & UML sequence diagram & Sequence Diagrams; Interface Analysis; Specify and Model Requirements.\\
How does an entity behave over its lifecycle? & UML state machine / state model & State Modelling; Business Rules Analysis; Validate Requirements.\\
How is the domain structured? & UML class/domain model where appropriate & Data Modelling; Concept Modelling; Data Dictionary.\\
How does business work flow across roles and systems? & BPMN collaboration/process model & Process Modelling; Process Analysis; Roles and Permissions Matrix.\\
Where is the system boundary and what crosses it? & Context/scope model plus interfaces; UML component/context views if useful & Scope Modelling; Interface Analysis; Data Flow Diagrams.\\
\bottomrule
\end{longtable}

\begin{warningbox}
\textbf{Notation discipline.} UML and BPMN answer different questions. BPMN is strongest for business-process flow, events, gateways, roles, and handoffs. UML is strongest for software-oriented structure, interactions, behavior, and actor-goal views. Use both only when each contributes a distinct viewpoint.
\end{warningbox}
\section{AppSec Integration Pattern}
\begin{enumerate}
\item \textbf{Strategy Analysis:} identify business/security need, current attack surface, critical assets, compliance context, target security outcomes, and risk tolerance.
\item \textbf{Planning and Monitoring:} define security stakeholders, risk/exception governance, repositories, traceability, and evidence expectations.
\item \textbf{Elicitation and Collaboration:} elicit abuse cases, trust boundaries, data classifications, security assumptions, external dependencies, and operational constraints.
\item \textbf{RADD:} specify security functional requirements and NFRs; model trust-sensitive flows; verify testability; validate risk reduction; compare control/design options.
\item \textbf{Requirements Life Cycle Management:} trace security requirements to threats, controls, code/components, automated checks, penetration-test findings, waivers, and releases.
\item \textbf{Solution Evaluation:} measure actual security outcomes and operational burden; analyze escaped vulnerabilities, control failures, remediation latency, and enterprise limitations.
\end{enumerate}
\section{Agile Integration Pattern}
An Agile interpretation does not eliminate analysis; it changes timing and granularity. Maintain enough upstream reasoning to preserve outcome and risk context, then defer detail until it can influence a near-term decision. Useful practices include:
\begin{itemize}
\item maintain outcome-oriented epics/features and trace them to smaller work items;
\item treat user stories as one representation, not as the only allowable requirement form;
\item attach acceptance/evaluation criteria and examples where they improve shared understanding;
\item refine continuously and re-prioritize as value, risk, evidence, or dependencies change;
\item use prototypes, spikes, and operational data to reduce uncertainty;
\item evaluate completed increments against actual outcomes, not only story completion.
\end{itemize}
\section{Enterprise Architecture Integration Pattern}
The strongest bridge from business analysis to enterprise architecture is the chain from \textbf{need and value} to \textbf{current/future state}, then to \textbf{capabilities and requirements}, then to \textbf{design options and transition states}, and finally to \textbf{measured outcomes}. Business Capability Analysis, Scope Modelling, Organizational Modelling, Process Modelling, Interface Analysis, and Requirements Architecture are especially useful for maintaining that chain.

\begin{practicebox}
\textbf{Architecture review question set.} What stakeholder concern is being addressed? Which business requirement and capability does it support? Which design decision satisfies it? What transition is required? What evidence will show that value was realized? What new risks or constraints does the decision introduce?
\end{practicebox}
\section{Technique-to-Domain Hot List}
\begin{longtable}{L{0.22\textwidth} L{0.66\textwidth}}
\toprule
Software engineering & Acceptance Criteria; Data Modelling; Interface Analysis; Functional Decomposition; Non-Functional Requirements Analysis; Prototyping; Reviews; Sequence Diagrams; State Modelling; Use Cases and Scenarios.\\
Application security & Risk Analysis and Management; Roles and Permissions Matrix; Non-Functional Requirements Analysis; Scope Modelling; Data Flow Diagrams; Interface Analysis; Sequence Diagrams; State Modelling; Vendor Assessment; Acceptance Criteria.\\
Agile & Backlog Management; User Stories; Acceptance Criteria; Prioritization; Prototyping; Workshops; Reviews; Estimation; Lessons Learned.\\
UML / BPMN & Process Modelling; Process Analysis; Sequence Diagrams; State Modelling; Use Cases and Scenarios; Data Modelling; Scope Modelling; Roles and Permissions Matrix.\\
Enterprise architecture & Business Capability Analysis; Business Cases; Business Model Canvas; Organizational Modelling; Scope Modelling; Process Modelling; Benchmarking; SWOT; Financial Analysis; Metrics and KPIs.\\
\bottomrule
\end{longtable}

\section*{Primary Source}
\addcontentsline{toc}{section}{Primary Source}
International Institute of Business Analysis (IIBA), \emph{A Guide to the Business Analysis Body of Knowledge (BABOK Guide)}, Version 3, 2015. This document is a study and practitioner synthesis of the user-supplied copy; it is not an official IIBA publication.
\end{document}
